\scanpage{135}
% archive foliation: 134
\deleted{\unsure{d)} Considérons donc la factorisation de Stein.}
\diagram{f\quad\left[\begin{tikzcd}X\arrow[r]&\unsure{Y'}\arrow[r,dash]&\unsure{Y}\end{tikzcd}\right]\unsure{\downarrow}}

On est ramené au cas où $Y$ est \uline{réduit}.

c) Prenons la factorisation de Stein
\diagram{\begin{tikzcd}X\arrow[r,dash]&Y'\arrow[r,dash]&\unsure{Y}\end{tikzcd}}
où $Y'\mathbin{\text{---}}Y$ est fini surjectif, donc un épimorphisme fini factorisable en produit d'épimorphismes stricts. Par ``\unsure{Sorb. chev.}'', $\uline{\mathrm{Pic}}_Y\mathbin{\text{---}}\uline{\mathrm{Pic}}_{Y'}$ est affine. On est ramené

\deleted{à prouver que $\uline{\mathrm{Pic}}_{Y'}\longrightarrow\uline{\mathrm{Pic}}_X$ l'est, \unsure{ce qui}}

\figure{\begin{tikzpicture}\draw[dashed] (0,0)--(12,0);\end{tikzpicture}}

ce qui nous ramène au cas $f_*(O_X)\simeq O_Y$.

d) Mais dans ce cas \unsure{on sait aussi que} $\uline{\mathrm{Pic}}_X\mathbin{\text{---}}\uline{\mathrm{Pic}}_Y$ est un mono, donc \unsure{de} \deleted{\unsure{noyau nul}} une immersion fermée, donc affine.

\margin{Attention au fait que $\uline{\mathrm{Pic}}_X$ et $\uline{\mathrm{Pic}}_Y$ ne sont pas de type fini sur $k$, \unsure{un mono n’est donc pas} néc\unsure{!} affine [ex. $\mathbf{Z}\to G$, $G=G_a$ ou bien VA]}

N.B. On s'est servi de l'existence des $\uline{\mathrm{Pic}}_Z$ pour $Z=Y_{\mathrm{réd}},X_{\mathrm{réd}}$, et pour les préschémas finis sur $Y_{\mathrm{réd}}$; \deleted{\unsure{Et}} on peut sans doute s'en passer !

